\documentclass[12pt]{letter}
\usepackage{geometry}
\geometry{margin=1in}
\usepackage{hyperref}

\begin{document}

\begin{letter}{%
Editor-in-Chief\\
\emph{Journal of Medical Ethics}\\
BMJ Publishing Group}

\opening{Dear Editor,}

I am pleased to submit the manuscript ``Geometric Clinical Ethics: Medical Decision-Making as Pathfinding on the Moral Manifold, and the Mathematical Structure of Moral Injury'' for consideration as an Extended Essay in the \emph{Journal of Medical Ethics}.

The paper provides the mathematical framework that clinical ethics has lacked.  It constructs a nine-dimensional clinical decision complex on which clinical decisions are pathfinding problems, models the clinician's decision as $A^*$ search with evidence-based medicine as the accumulated cost and clinical wisdom as the heuristic, and identifies the minimum-cost path---the clinical geodesic---as the replacement for scalar QALY optimization.

The paper makes six principal contributions:
\begin{enumerate}
    \item \textbf{Principlism as geometry:} Beauchamp and Childress's four principles are dimensions of the clinical manifold; conflicts between them are resolved by the metric, not by ad hoc balancing.
    \item \textbf{QALY irrecoverability:} A formal proof that scalar projection destroys irrecoverable moral information, with a precise characterization of when QALY analysis is and is not adequate.
    \item \textbf{Moral injury as forced boundary crossing:} The first mathematical characterization of moral injury, structurally distinct from burnout, with a predictive Moral Injury Index that enables prospective evaluation of institutional policies.
    \item \textbf{Informed consent as gauge invariance:} A functional test of consent validity---does the patient's decision survive reframing?---that goes beyond the legal checklist.
    \item \textbf{Geometric triage:} Triage as constrained multi-agent pathfinding, with a theorem predicting when utilitarian and manifold-optimal triage diverge.
    \item \textbf{The clinical Bond Index:} A quantitative measure of structural healthcare injustice, enabling evidence-based health equity policy.
\end{enumerate}

The framework generates six falsifiable predictions about clinician moral distress, consent quality, triage resistance, and health equity that differ from predictions of standard bioethics and QALY-based health economics.

This work is relevant to the \emph{Journal of Medical Ethics} because it provides a unified mathematical foundation for clinical ethical reasoning that connects principlism, moral injury research, informed consent theory, triage ethics, and health equity measurement---domains that have been addressed separately in the clinical ethics literature.

The manuscript has not been published previously and is not under consideration elsewhere.  No human participants were involved; this is a theoretical paper.  The parent framework is publicly available under the MIT open-source licence.

Generative AI tools (Claude, Anthropic) were used for manuscript preparation and formatting; all intellectual content is the author's own.

Thank you for considering this submission.

\closing{Sincerely,\\[2em]
Andrew H. Bond\\
Senior Lecturer, Department of Computer Engineering\\
San Jos\'{e} State University\\
\href{mailto:andrew.bond@sjsu.edu}{andrew.bond@sjsu.edu}\\
ORCID: 0009-0003-2599-6158}

\end{letter}
\end{document}
